\begin{statement}{G.22}{Proved}{Prym detection and the condition for punctured elliptic transport}
Let $k/K$ be finite and let $p:Y\to U_k$ be a geometrically
connected finite \emph{\'etale} cover of degree $d$. Write
$\overline p:\overline Y\to E_k$ for its extension to smooth
proper curves, $S=\overline Y\setminus Y$, and
$g_Y=g(\overline Y)$; cardinalities of boundary sets are geometric.
The following assertions delimit the finite-cover attack on G.20.
\begin{enumerate}
\item If $u\in\pi_1(Y_{\overline k})$ has image in $R$, its
proper homology class $\overline v_\ell$ lies in
\[
 \ker\bigl(V_\ell J_{\overline Y}
              \xrightarrow{\overline p_*}V_\ell E\bigr)
       =V_\ell\bigl(\ker(\overline p_*:J_{\overline Y}\to E)^0\bigr).
\]
In particular $\overline v_\ell=0$ if $g_Y=1$. A residual
weight-one detector must therefore occur in the Prym part of a
cover of genus at least two.
\item For a nonconstant $f:\overline Y\to E'$ and a geometric
point $q\in E'$, there is a restriction
$Y\to E'\setminus\{q\}$ exactly when
\[
                    \operatorname{Supp}f^{-1}(q)\subset S.
\]
This condition implies
\[
                 \deg f\le |S|+2g_Y-2=d.
\]
Equality holds exactly when $S=\operatorname{Supp}f^{-1}(q)$
and $f$ is unramified away from that fibre; then the restricted
map is finite \emph{\'etale}.
\item Suppose $f$ has this fibre condition and is defined, with
$q$, over a finite extension of $k$. Suppose also that the
pro-$\ell$ arithmetic normalizer criterion of G.13 holds for
the \emph{particular pointed curve} $(E',q)$. No procyclic
subgroup of $\pi_1(Y_{\overline k})$ with open arithmetic
normalizer can have proper homology image carried nontrivially
by $f_*$ to $V_\ell E'$. The target hypothesis is established
here when $(E',q)$ becomes isomorphic to $(E,O)$ after a finite
extension. It is not asserted for arbitrary isogenous curves,
CM fields, or higher-dimensional CM factors.
\item A detecting elliptic Jacobian quotient need not satisfy
the fibre condition. Explicitly, $[2]:E\setminus E[2]\to U$
is a degree-four finite \emph{\'etale} cover, and
$[3]:E\to E$ is injective on every Tate module, but
\[
                      [3](E\setminus E[2])=E.
\]
It therefore gives no map to any once-punctured $E$.
\item Adding the missing punctures does not in general permit
an arithmetic-normalizer-preserving lift of a procyclic subgroup.
A primitive CM eigenline in $T_5E$ has open arithmetic normalizer,
but no procyclic subgroup of $P_5$ with open arithmetic normalizer
maps nontrivially into that line under the filling homomorphism
$P_5\twoheadrightarrow T_5E$.
\end{enumerate}
\end{statement}
\begin{proof}
The map $\pi_1(Y_{\overline k})\to\pi$ is the geometric
inclusion determined by a basepoint path. The induced homomorphism
to the abelian pro-$\ell$ quotient $T_\ell E$ factors through
$T_\ell J_{\overline Y}$ and $\overline p_*$. It vanishes on
an element mapping into $R$, because this quotient factors
through $P_\ell$. Tensoring with $\mathbb Q_\ell$ proves the
kernel assertion. Factoring
the surjective homomorphism $J_{\overline Y}\to E$ by the
identity component of its kernel leaves an isogeny, which is
injective on Tate modules; this gives the displayed equality.
If $g_Y=1$, Riemann--Hurwitz makes $\overline p$ unramified.
After a finite extension and choices of origins it is an isogeny
of elliptic curves. Its Tate-module map is injective: a compatible
torsion sequence in its finite kernel must be zero, since a
fixed power of $\ell$ kills that kernel's $\ell$-part. This
proves the genus-one assertion. In particular the original
elliptic quotient of $J_{\overline Y}$ cannot itself carry the
residual vector; any relevant CM component must occur in the
indicated kernel.

The fibre condition is just the condition that the restriction
avoid $q$. All ramification of $\overline p$ lies above $O$,
and that fibre is supported on $S$. In characteristic zero,
Riemann--Hurwitz gives
\[
 2g_Y-2=\sum_{P\in S}(e_P(\overline p)-1)=d-|S|.
\]
For $f$ it gives $\sum_P(e_P(f)-1)=2g_Y-2$; see
\citet[\S53.12, Lemmas 53.12.2 and 53.12.4]{GStacksRiemannHurwitz}.
Under the fibre condition,
\[
 \deg f
   =|f^{-1}(q)|+\sum_{P\in f^{-1}(q)}(e_P(f)-1)
   \le |S|+2g_Y-2=d.
\]
Here the first cardinality counts distinct geometric points.
Equality requires equality in both inequalities, precisely as
stated; the resulting restriction is finite and unramified
between smooth curves, hence finite \emph{\'etale}.

For the third assertion, make the finite base extension on which
all maps are defined. An open arithmetic image stays open under
intersection with the corresponding open arithmetic fundamental
group. Choose compatible geometric basepoints and a path. The
morphism $Y\to E'\setminus\{q\}$ induces a homomorphism of
arithmetic fundamental groups over $G_k$, and hence one to the
geometrically pro-$\ell$ target. It carries a normalizer into
the normalizer of the image subgroup and preserves its open
arithmetic image. The nonzero proper homology image ensures
that the image subgroup is a nontrivial closed procyclic
subgroup of the target geometric group. The assumed normalizer
criterion puts it in cusp inertia. But cusp inertia has zero
image in $T_\ell E'$, contradicting the detecting hypothesis.
When the pointed target becomes $(E,O)$, this is exactly G.13:
the finite extension changes the arithmetic group to an open
subgroup, and changes of basepoint only conjugate the relevant
subgroups. No claim about a different CM curve follows from
this argument merely because its Jacobian has an abelian
Galois representation.

For application to an original normalizer, one can first choose
an open geometric cover containing its procyclic subgroup.
Intersecting its normalizer with an arithmetic fundamental
group of that descended cover leaves an open arithmetic image,
so the preceding functorial argument applies. The missing step
is a suitable $f$ satisfying both detection and the fibre
condition. A CM character on a Tate module alone supplies neither
a rational elliptic quotient nor this condition: a CM factor
can have dimension greater than one, and its individual
$\mathbb Q_\ell$ eigenlines need not come from elliptic
quotients.

For the explicit example, $E[2]=\{O,(0,0),(1,0),(-1,0)\}$,
and $[2]$ is unramified in characteristic zero. Every fibre of
$[3]$ has nine geometric points. Since $[3]$ acts as the
identity on $E[2]$, a fibre over $q\in E[2]$ contains exactly
one point of $E[2]$ and eight points of its complement. A fibre
over $q\notin E[2]$ is wholly in the complement. This proves
the displayed surjectivity. The proper homology map is $3I$,
which is injective even over $\mathbb Z_3$.
The exact independent certificate in
\texttt{code/g\_cover\_transport.py} derives
\[
 \psi_3(x)=3x^4-6x^2-1,
 \qquad x(2P)-x(P)=\frac{-\psi_3(x(P))}{4(x(P)^3-x(P))}.
\]
It verifies that $\psi_3$ is squarefree and relatively prime
to $x^3-x$, and that its values at $0,1,-1$ are $-1,-4,-4$.
Thus its four roots give eight nonzero points of order three,
all outside $E[2]$, in agreement with the fibre argument.
This example refutes only automatic puncturing of a detecting
quotient; it is not a residual element or a counterexample to
G.20. More generally, postcomposing any detecting $f$ by $[m]$
preserves nonzero Tate-module detection but multiplies its
degree by $m^2$. For large $m$ the degree bound excludes every
puncture. Passing to a further surjective cover of $Y$ cannot
repair a fixed map whose image is all of $E'$.

Finally, CM over $K$ and the splitting of $5$ give a primitive
$\mathbb Z_5$ eigenline $L\subset T_5E$ stable under $G_K$.
The zero section splits the arithmetic fundamental group of the
proper elliptic curve, so the normalizer of $L$ has full
arithmetic image. If a closed procyclic subgroup upstairs in
$P_5$ mapped nontrivially to $L$ and retained an open arithmetic
normalizer, G.13 would make it cuspidal. Filling the puncture
kills its image, a contradiction. This failure persists after
any finite base extension. Likewise G.10 shows that a nonzero
$\Delta$-fixed vector of $T_5E$ has no lift fixed by any of
the specified cusp-packet actions on $P_5$. Thus deleting more
points cannot be used as an automatic equivariant lifting step.
The conditional transport excludes a residual normalizer
\emph{when} the required map exists; it does not establish that
existence for the covers in G.20.
\end{proof}
